\documentclass[a4paper,12pt]{article}
\usepackage[utf8]{inputenc}
\usepackage[T1]{fontenc}
\usepackage[french]{babel}
\usepackage{geometry}
\geometry{left=2.5cm,right=2.5cm,top=2.5cm,bottom=2.5cm}

\usepackage{lmodern}
\usepackage{xcolor}
\definecolor{titleblue}{RGB}{0,51,140}
\definecolor{TITLEBLUE}{RGB}{0,51,140}

\usepackage{hyperref}

\pagestyle{empty}

\begin{document}

\begin{flushleft}
    \textbf{Ismail AISSAOUI} \\
    Ingénieur d'État en Intelligence Artificielle \\
    Chlef, Algérie \\
    ismail.aissaoui.pro@gmail.com \\
    +213 660 70 77 96 \\
    \href{https://ismail-aissaoui.vercel.app}{ismail-aissaoui.vercel.app}
\end{flushleft}

\begin{flushright}
    À l'attention du \\
    \textbf{Professeur Yacine OUZROUT} \\
    Laboratoire DISP \\
    Université Lumière Lyon 2 \\
    69621 Villeurbanne, France \\
    \medskip
    Le 29 Avril 2026
\end{flushright}

\bigskip

\noindent \textbf{Objet : Candidature pour un doctorat en Maintenance Prédictive \& IA Multimodale}

\bigskip

Monsieur le Professeur,

C'est avec un grand intérêt que je vous adresse ma candidature pour poursuivre des travaux de doctorat sous votre direction au sein du laboratoire DISP. Vos travaux sur la maintenance prédictive pour les systèmes industriels complexes et l'intégration de l'IA multimodale résonnent fortement avec mon parcours d'ingénieur.

Actuellement stagiaire à la Sonatrach, je travaille sur la mise en œuvre d'un \textbf{Jumeau Numérique (Architecture Générative Mamba-SSM)} pour turbines à gaz. Ma recherche se concentre sur l'estimation de la durée de vie résiduelle (RUL) en utilisant une architecture duale combinant un modèle "Edge" (\textbf{C++ JIT xLSTM}) et une approche de substitut génératif pour optimiser la précision du diagnostic sous contrainte matérielle.

Je suis particulièrement intéressé par la manière dont nous pouvons coupler les modèles génératifs et les systèmes de recherche d'information pour assister les experts en maintenance dans le diagnostic de pannes critiques. Votre approche holistique de la performance des systèmes de production est le cadre idéal pour développer ces outils de soutien à la décision.

Je joins mon CV à cette lettre, en espérant avoir l'opportunité de discuter prochainement de projets de recherche potentiels pour 2026.

Je vous prie d'agréer, Monsieur le Professeur, l'expression de mes salutations les plus distinguées.

\bigskip

\begin{flushright}
    \textbf{Ismail AISSAOUI}
\end{flushright}

\end{document}
